\subsection{The Dark Dimension and the Swampland --- arXiv:2205.12293}

\textbf{Authors:} Miguel Montero, Cumrun Vafa, Irene Valenzuela

\paragraph{主な主張}
Swampland の Distance/Duality conjecture と観測された小さな dark-energy scale を組み合わせると、$m_{\rm tower}\sim \Lambda^{1/4}$ の light tower が必要となり、既存の重力・天体物理制限を踏まえると長さ $l\sim \Lambda^{-1/4}\sim 10^{-6}\,\mathrm{m}$ の単一の mesoscopic extra dimension に導かれると主張する \cite{Montero:2022prj}。さらに、その bulk を伝播する massless fermion は sterile-neutrino KK tower の候補となり、active neutrino mass scale $m_\nu\sim \langle H\rangle^2\Lambda^{-1/12}M_{\rm Pl}^{-2/3}$ を与える。また higher-dimensional Planck scale に対応する species scale を $\widehat M\sim \Lambda^{1/12}M_{\rm Pl}^{2/3}\sim 10^{9}$--$10^{10}\,\mathrm{GeV}$ と予測する。

\paragraph{新規性}
dark-energy scale を単なる低エネルギーの宇宙論パラメータとしてではなく Swampland の light-tower 条件と結び付け、micron-size extra dimension、sterile-neutrino tower、neutrino mass、species scale までを一つのスケーリング関係から関連付けた点が特徴的である。

\paragraph{位置付け}
large-extra-dimension neutrino model に量子重力側から新しい動機を与える代表的な ``dark dimension'' 論文であり、bulk right-handed neutrino / sterile KK tower を用いる現象論と Swampland を直接接続する。特に $S^1$ 型 bulk sterile-neutrino model の理論的動機として重要である一方、重力も同じ大きな余剰次元を感じるという前提では単一の大きな次元を強く選好するため、$T^2$ 型拡張を考える際には対照となる基準点である。
